%! Author = markt
%! Date = 3/10/2025

% Preamble
\documentclass[11pt]{article}

% Packages
\usepackage{latexsym,amsfonts,amssymb,amsthm,amsmath, yfonts, cancel, listings, xcolor, graphicx}
\usepackage{color}
\usepackage{listings}

\author{Mark Stanley}

\definecolor{codegreen}{rgb}{0,0.6,0}
\definecolor{codegray}{rgb}{0.5,0.5,0.5}
\definecolor{codepurple}{rgb}{0.58,0,0.82}
\definecolor{backcolour}{rgb}{0.95,0.95,0.92}
\lstdefinestyle{mystyle}{
    backgroundcolor=\color{backcolour},
    commentstyle=\color{codegreen},
    keywordstyle=\color{magenta},
    numberstyle=\tiny\color{codegray},
    stringstyle=\color{codepurple},
    basicstyle=\ttfamily\footnotesize,
    breakatwhitespace=false,
    breaklines=true,
    captionpos=b,
    keepspaces=true,
    numbers=left,
    numbersep=5pt,
    showspaces=false,
    showstringspaces=false,
    showtabs=false,
    tabsize=2
}

\title{Math 443 Session 26}

% Document
\begin{document}
    \lstset{style=mystyle}

    \maketitle

    \subsection{Characteristic Polynomial}
    Given some matrix $A$, the characteristic polynomial is defined as the solution to
    \[\det(A-\lambda I)=0\]
    Where $\lambda$ is the eigenvalue of the matrix $A$, and $I$ is the identity matrix.
    The solution to this is a polynomial.

    \subsection{Leading Coefficients}
    From expansion, we have:
    \[p_A(\lambda) = \det(A-\lambda I)=c_n \lambda^n + c_{n-1}\lambda^{n-1}+\cdots+c_1\lambda + c_0\]
    We can identify $c_0$ is a constant, $c_n=(-1)^n$, $c_{n-1}=(-1)^{n-1} tr(A)$, where
    \[tr(A)=a_{11}+a_{22} + \cdots + a_{nn}\]

    \subsection{Example}
    Given some $2 \times 2$ matrix $A=\begin{pmatrix} a & b \\ c & d \end{pmatrix}$, we can find the characteristic polynomial as follows:
    \[p_A(\lambda)=\lambda^2 - (tr(A))\lambda + \det(A)\]

    \theorem{Fundamnetal Theorem of Algebra}
    The fundamental theorem of algebra allows us to factor the characteristic polynomial into linear factors.
    \[p_A(\lambda)=(-1)^n(\lambda-\lambda_1)(\lambda-\lambda_2)\ldots(\lambda-\lambda_n)\]
    Where $\lambda_i$ are the eigenvalues of the matrix $A$.

    \theorem{8.11}
    An $n \times n$ matrix possesses at least one and at most $n$ distinct complex eigenvalues

    Some code:
    \begin{lstlisting}[language=Python]
    def python_code():
        pass
    \end{lstlisting}

\end{document}